\section{Maximal estimates and principal values}\label{sec:consequences}

The counting estimate already contains the maximal estimate, by taking
a single annulus. We first record its continuity with respect to the
inputs, then use that continuity to transfer joint principal-value
convergence from smooth functions to all complex $L^3$ inputs.
Throughout this section, write $p=3/2$.

\begin{corollary}[The hard maximal operator]\label{cor:maximal}
For complex $F,G\in L^3(\R^2)$, define
\[
 B_*(F,G)(x,y)
 =\sup_{0<\varepsilon<R<\infty}|B_{\varepsilon,R}(F,G)(x,y)|
\]
on the common full-measure set of Lemma~\ref{lem:finite-truncations},
and set it to zero elsewhere. This function is measurable, and there is
an absolute constant $C_M$ such that
\begin{equation}\label{eq:hard-maximal-bound}
 \|B_*(F,G)\|_p\le C_M\|F\|_3\|G\|_3.
\end{equation}
The supremum includes both finite hard truncation endpoints. Moreover,
$B_*$ is a continuous map from $L^3(\R^2)\times L^3(\R^2)$ to
$L^p(\R^2)$ and is the unique continuous extension of its values on
compactly supported smooth pairs.
\end{corollary}

\begin{proof}
Endpoint continuity in Lemma~\ref{lem:finite-truncations} identifies the
supremum with its restriction to rational endpoint pairs. Consequently
it is measurable and equals $\mathcal S_1(F,G)$. Theorem~\ref{thm:count}
gives \eqref{eq:hard-maximal-bound}, with its absolute factor
$\log 3$ absorbed into $C_M$.

For two pairs of inputs, bilinearity of each finite truncation gives,
almost everywhere,
\begin{equation}\label{eq:maximal-continuity}
 |B_*(F,G)-B_*(f,g)|
 \le B_*(F-f,G)+B_*(f,G-g).
\end{equation}
Taking $L^p$ norms and using \eqref{eq:hard-maximal-bound} proves
continuity. Density of compactly supported smooth functions gives the
asserted uniqueness.
\end{proof}

\begin{corollary}[Joint bilinear principal values]\label{cor:bilinear-pv}
For every complex $F,G\in L^3(\R^2)$, the joint limit
\[
 B(F,G)=\lim_{\substack{\varepsilon\downarrow0\\R\uparrow\infty}}
 B_{\varepsilon,R}(F,G)
\]
exists almost everywhere and in $L^p(\R^2)$, and
\[
 \|B(F,G)\|_p\le C_M\|F\|_3\|G\|_3.
\]
More precisely, the entire tail converges in the maximal sense
\begin{equation}\label{eq:maximal-tail-convergence}
 \left\|\sup_{\substack{0<\varepsilon<1/n\\R>n}}
 |B_{\varepsilon,R}(F,G)-B(F,G)|\right\|_p
 \longrightarrow0.
\end{equation}
The operator $B$ is the unique bounded complex bilinear extension of
the principal-value operator on compactly supported smooth inputs.
\end{corollary}

\begin{proof}
For compactly supported smooth $f,g$ and fixed $(x,y)$, put
$h(t)=f(x+t,y)g(x,y+t)$. Then
\[
 B_{\varepsilon,R}(f,g)(x,y)
 =\int_\varepsilon^R\frac{h(t)-h(-t)}t\,\dd t.
\]
The numerator is $O(t)$ at zero and vanishes for large $t$, so the joint
limit exists at every point.

For general $F,G$, define the tail oscillation
\[
 \Omega_n(F,G)=
 \sup_{\substack{0<\varepsilon,\varepsilon'<1/n\\R,R'>n}}
 |B_{\varepsilon,R}(F,G)-B_{\varepsilon',R'}(F,G)|.
\]
These functions are measurable by endpoint continuity, decrease with
$n$, and satisfy $0\le\Omega_n(F,G)\le2B_*(F,G)$. Choose compactly
supported smooth $f_m\to F$ and $g_m\to G$ in $L^3$. Take a common
full-measure set for the pairs $(F,G)$, $(f_m,g_m)$, $(F-f_m,G)$,
and $(f_m,G-g_m)$, for all $m$. On this set, bilinearity gives
\[
 \Omega_n(F,G)\le\Omega_n(f_m,g_m)
       +2B_*(F-f_m,G)+2B_*(f_m,G-g_m).
\]
The smooth joint convergence implies $\Omega_n(f_m,g_m)\to0$ pointwise.
Thus, writing $\Omega_\infty=\lim_n\Omega_n(F,G)$, we obtain
\[
 \|\Omega_\infty\|_p
 \le2C_M\bigl(\|F-f_m\|_3\|G\|_3
             +\|f_m\|_3\|G-g_m\|_3\bigr).
\]
Letting $m\to\infty$ shows that $\Omega_\infty=0$ almost everywhere.
This is the joint Cauchy criterion for the two endpoints, and hence
defines $B(F,G)$ almost everywhere. Taking, for example, the cofinal
sequence $(\varepsilon,R)=(1/(2n),2n)$ shows that the limit is measurable;
its magnitude is at most $B_*(F,G)$. Define it to be zero on the
remaining null set.

Letting $(\varepsilon',R')$ tend jointly to $(0,\infty)$ in the
definition of $\Omega_n$ gives
\[
 \sup_{\substack{0<\varepsilon<1/n\\R>n}}
 |B_{\varepsilon,R}(F,G)-B(F,G)|\le\Omega_n(F,G).
\]
The supremum is measurable, again by endpoint continuity.
Dominated convergence applies because
$\Omega_n(F,G)\le2B_*(F,G)\in L^p$, proving
\eqref{eq:maximal-tail-convergence} and therefore joint norm convergence.
The norm bound follows from $|B(F,G)|\le B_*(F,G)$. Passing to the
$L^p$ limit in the bilinear identities for finite truncations proves
complex bilinearity. The bound and density then prove uniqueness.
\end{proof}

\subsection{Flat and simplex scalar forms}

For complex $F_0,F_1,F_2\in L^3(\R^2)$, let
\[
 \mathcal L_{\varepsilon,R}(F_0,F_1,F_2)
 =\iiint_{\varepsilon<|t|<R}
 F_0(x,y)F_1(x+t,y)F_2(x,y+t)
 \,\dd x\,\dd y\,\frac{\dd t}{t}.
\]
No input is conjugated in this scalar form.

\begin{corollary}[Flat scalar principal values]\label{cor:scalar-pv}
For every complex $L^3$ triple, each finite flat truncation is absolutely
integrable, and
\[
 \sup_{0<\varepsilon<R<\infty}
 |\mathcal L_{\varepsilon,R}(F_0,F_1,F_2)|
 \le C_M\prod_{v=0}^2\|F_v\|_3.
\]
Its joint scalar principal value exists and equals
$\int_{\R^2}F_0 B(F_1,F_2)$. This limiting form is the unique bounded
complex trilinear extension of its compactly supported smooth definition.
\end{corollary}

\begin{proof}
H\"older's inequality in $(x,y)$ bounds the integral of the absolute
value by $2\log(R/\varepsilon)\prod_v\|F_v\|_3$. Fubini's Theorem
therefore identifies the finite form with
$\int F_0 B_{\varepsilon,R}(F_1,F_2)$.
Corollary~\ref{cor:maximal} gives its uniform bound, and pairing the
joint $L^p$ limit in Corollary~\ref{cor:bilinear-pv} with $F_0$ gives
the scalar principal value. The limiting form is bounded and complex
trilinear; density proves uniqueness.
\end{proof}

In simplex coordinates the corresponding scalar form is
\[
 \Lambda_{\varepsilon,R}(G_0,G_1,G_2)
 =\iiint_{\varepsilon<|a+b+c|<R}
 G_0(a,b)G_1(b,c)G_2(c,a)
 \,\frac{\dd a\,\dd b\,\dd c}{a+b+c}.
\]
The equivalence of the flat and simplex formulations, including their
optimal constants, is described in
\cite[Appendix~B.1]{KovacThieleZorinKranich2015}. We record the
norm-preserving maps to track both hard endpoints and their joint limit.

\begin{proposition}[Scalar equivalence]\label{prop:flat-equivalence}
For every complex $G_0,G_1,G_2\in L^3(\R^2)$, each finite simplex
truncation is absolutely integrable, and
\[
 \sup_{0<\varepsilon<R<\infty}
 |\Lambda_{\varepsilon,R}(G_0,G_1,G_2)|
 \le C_M\prod_{v=0}^2\|G_v\|_3.
\]
Its joint scalar principal value exists and is the unique bounded
complex trilinear extension of its compactly supported smooth
definition. The optimal constants in the uniform finite flat and
simplex scalar estimates are equal.
\end{proposition}

\begin{proof}
Use the determinant-one change of variables
\[
 x=-a,\qquad y=-b,\qquad t=a+b+c,
 \qquad
 a=-x,\quad b=-y,\quad c=t+x+y.
\]
For flat inputs set
\[
 \begin{aligned}
 G_0(u,v)&=F_0(-u,-v),\\
 G_1(u,v)&=F_1(u+v,-u),\\
 G_2(u,v)&=F_2(-v,u+v).
 \end{aligned}
\]
Each pair map has determinant one. The inverse formulas are
\[
 \begin{aligned}
 F_0(r,s)&=G_0(-r,-s),\\
 F_1(r,s)&=G_1(-s,r+s),\\
 F_2(r,s)&=G_2(r+s,-r).
 \end{aligned}
\]
These maps preserve $L^3$ norms, compact smoothness, and the Schwartz
class. Direct substitution identifies both the integrands and the
finite truncation regions. Absolute integrability, the uniform bound,
and joint scalar convergence therefore follow from
Corollary~\ref{cor:scalar-pv}. The inverse maps realize every simplex
triple, so the optimal constants are equal. Density proves uniqueness
of the bounded trilinear extension.
\end{proof}

For completeness, let $C_S$ denote this common optimal scalar constant.
It is also the optimal constant for the uniform family of individual
bilinear bounds
\[
 \|B_{\varepsilon,R}(F,G)\|_p
 \le C_S\|F\|_3\|G\|_3.
\]
Indeed, H\"older's inequality gives one direction. For the other, if
$U\in L^p$ is nonzero, then
\[
 H=\frac{\overline U}{|U|}\,
       \frac{|U|^{1/2}}{\|U\|_p^{1/2}},
\]
defined to be zero where $U=0$, satisfies $\|H\|_3=1$ and
$\int HU=\|U\|_p$. Apply this duality formula to each finite bilinear
truncation. In particular, the optimal pointwise maximal constant is
at least $C_S$; equality of the scalar constants does not assert
equality with the maximal or variation constants.
